\documentclass[10pt]{article}
\usepackage[a4paper,margin=1.85cm]{geometry}
\usepackage{amsmath,amssymb,amsthm,booktabs,enumitem}
\usepackage{xcolor}
\usepackage[colorlinks=true,linkcolor=blue!55!black]{hyperref}

\theoremstyle{plain}
\newtheorem{theorem}{Theorem}
\newtheorem{proposition}{Proposition}

% one-line provenance tag under an equation
\newcommand{\src}[1]{\par\smallskip\noindent{\small\emph{From:} #1}\par\smallskip}

\setlist[itemize]{itemsep=1pt,topsep=3pt,leftmargin=1.2em}
\setlength{\parskip}{2pt}
\linespread{0.97}

\title{\vspace{-1.4cm}\bfseries The Model in Three Pages\\
\large Multinational ownership, market power, and the measurement of foreign profits}
\date{\vspace{-0.3cm}\small\today\ --- working draft, please do not circulate}

\sloppy
\raggedbottom
\begin{document}
\maketitle
\vspace{-0.7cm}

\section{Introduction}

A large share of Latin American exports is shipped by subsidiaries of international
groups headquartered elsewhere --- often in the very country imposing the tariff.
\textbf{How much of the profit earned in those markets belongs to that country?}
The obvious answer --- total profit times the share of firms owned --- is wrong, and
Proposition~\ref{p:main} gives the exact error. Getting there needs three things at
once: markups that rise with firm size, multinational production, and entry. Each is
standard; the combination, plus ownership recorded \emph{firm by firm}, is not.

\section{Model}

A \emph{market} is one product in one country. Factories $i$ sell into it; each
belongs to a \emph{group} $g$ that may own factories in several countries. $S_g$ is
the group's share of market sales; $\sigma>\eta\ge1$ are the elasticities of
substitution within and across markets.

\paragraph{1. The markup.} A group expanding one factory also steals sales from its
\emph{own} other factories, and it counts that loss. So the bigger the group, the more
it holds back:
\begin{equation}
\underbrace{\frac{1}{\mu(S)}=1-\frac{1-S}{\sigma}-\frac{S}{\eta}}_{\text{quantity competition}}
\qquad\qquad
\underbrace{\mu(S)=\frac{\sigma-(\sigma-\eta)S}{\sigma-(\sigma-\eta)S-1}}_{\text{price competition}}
\label{e:mu}
\end{equation}
Both rise from $\sigma/(\sigma-1)$ at $S=0$ toward $\eta/(\eta-1)$ at $S=1$; every
result below holds for \emph{any} increasing $\mu$, so how firms compete is a
parameter, not a commitment.
\src{Atkeson and Burstein (2008), with the markup moved from the factory to the
\emph{owner} --- the internalisation Yang (2023) calls cannibalisation.}

\paragraph{2. Profit.} With margin $L(S)\equiv1-1/\mu(S)$ and market spending $E$:
\begin{equation}
\Pi_g=E\,S_g\,L(S_g)
\;\;=\;\;\frac{E}{\sigma}\,S_g\bigl(1+\kappa S_g\bigr)\ \text{under quantities},
\qquad \kappa\equiv\frac{\sigma}{\eta}-1 .
\label{e:profit}
\end{equation}
Profit rises \emph{faster} than size; the $S^2$ term is the entire content of the
paper. Conduct only scales it ($\kappa=4.0$ quantities vs $0.8$ prices at
$\sigma=5,\eta=1$).
\src{The closed form is exact under Cournot with $\eta=1$; the linear margin
$L=(1+\kappa S)/\sigma$ is what makes Proposition~\ref{p:main} an identity.}

\paragraph{3. One market, one unknown.} With $K_g\equiv\sum_{i\in g}c_i^{1-\sigma}$
(the group's \emph{strength}, $c_i$ = delivered cost), $A\equiv\sum_j p_j^{1-\sigma}$
(how \emph{tough} the market is), and $\psi(S)\equiv S\,\mu(S)^{\sigma-1}$:
\begin{equation}
\boxed{\;S_g=\psi^{-1}\!\left(K_g/A\right),\qquad \textstyle\sum_g S_g=1.\;}
\label{e:share}
\end{equation}
However many factories a group owns, its whole position is one number, $K_g$ ---
cannibalisation is already inside it --- and all it needs to know about rivals is
$A$. A market with hundreds of factories is a problem in one unknown; that is what
makes the theorems possible.
\src{This model (the reduction is what the uniqueness proofs stand on).}

\paragraph{4. Costs.} Factory $a$ (head office $h$, located in $\ell$, selling to $n$,
sector $k$):
\begin{equation}
c_{a,n}=\frac{1}{\varphi_a}\,
\underbrace{w_\ell^{1-\nu_k}}_{\text{wages}}
\underbrace{\bigl(P^{\text{IO}}_{\ell k}\bigr)^{\nu_k}}_{\text{inputs}}
\underbrace{\gamma_{h\ell}}_{\text{away from HQ}}
\underbrace{d_{\ell n}}_{\text{shipping from }\ell}
\underbrace{(1+t_{\ell nk})}_{\text{tariff}}
\label{e:cost}
\end{equation}
Labour and inputs are bought where the factory sits; shipping is counted from the
factory, so one plant serves many countries. Being a multinational shows up in
$\varphi_a$: head-office services are a non-rival \emph{capability} (complexity
$\alpha_k$ penalises stand-alone locals and rewards group capability), never a foreign
wage inside a domestic cost --- the choice that decides whether the model has one
solution, and improves the fit besides.
\src{$\gamma_{h\ell}$: Ramondo--Rodr\'iguez-Clare (2013), measured by Head--Mayer
(2019). Shipping from the factory and the fixed cost below: Tintelnot (2017).
Inputs: Caliendo--Parro (2015). $\alpha_k$: Antr\`as (2003), Antr\`as--Helpman (2004),
used on the capability side.}

\paragraph{5. Entry.} Serving market $n$ costs $F_{a,n}=f_k\,w_\ell$ per period, paid
in the same labour as production --- so entry costs are real resources, they appear in
labour demand, and distributed profits are net of them. Which factories exist is an
outcome, not a list.
\src{Tintelnot (2017); the granular pool of potential parents, growing with country
size, is Gaubert--Itskhoki (2021).}

\paragraph{6. Income.} With $\theta_{gn}$ the share of group $g$ owned by country $n$:
\begin{equation}
X_n=\underbrace{w_nL_n}_{\text{wages}}
+\underbrace{\textstyle\sum_g\theta_{gn}\bigl(\Pi_g-\text{entry costs}_g\bigr)}_{\text{profits owned, net}}
+\underbrace{T_n}_{\text{tariffs}}
\label{e:income}
\end{equation}
Profits are earned where factories sell but \emph{received} where owners live; and
$\theta$ varies firm by firm rather than being one number per country.
\src{The firm-level $\theta$ is this project's own ingredient --- matched
Orbis/D\&B--customs records for ten Latin American origins; nothing in the prior
literature carries it.}

\paragraph{7. The answer.} For a country $H$ owning shares $\theta_g$, with
$\bar\theta=\sum_g S_g\theta_g$ and $\mathrm{Cov}(\theta,S)=\sum_g\theta_gS_g^2
-\bar\theta\,\mathrm{HHI}$:

\begin{proposition}[The correct adjustment is a concentration measure]\label{p:main}
\begin{equation}
\frac{\Pi_H}{E}=\frac1\sigma\Bigl[\bar\theta+\kappa\bigl(\bar\theta\,\mathrm{HHI}
+\mathrm{Cov}(\theta,S)\bigr)\Bigr],
\qquad
\boxed{\ \text{error of the country-level calculation}
=\tfrac{E\kappa}{\sigma}\,\mathrm{Cov}(\theta,S).\ }
\label{e:main}
\end{equation}
\end{proposition}

\noindent\emph{Proof.} Insert $L=(1+\kappa S)/\sigma$ into \eqref{e:profit}, sum with
weights $\theta_g$, and subtract $\bar\theta$ times total profit
$\Pi/E=(1+\kappa\,\mathrm{HHI})/\sigma$. $\square$

\noindent What matters is not what fraction of \emph{firms} a country owns but what
fraction of the \emph{big} ones. With constant markups ($\kappa=0$) the correction
vanishes --- which is why the standard workhorse models cannot host the question.
\textbf{Measured in our data} (market = destination$\times$product$\times$year, group
= ultimate parent): the country-level calculation understates the United States'
stake by $10\%$, overstates Argentina's and Peru's by $24$--$30\%$ --- the sign varies
owner by owner, which no country-level share can express.

\section{Existence and Uniqueness}

\begin{theorem}[One market]\label{t:market}
Given the operating factories ($\ge2$ groups), the market equilibrium exists and is
unique: $\psi$ is increasing, so $\sum_g\psi^{-1}(K_g/A)$ falls strictly in $A$ and
crosses $1$ exactly once; group profit is strictly concave in the group's own
quantities.
\end{theorem}

\begin{theorem}[Entry]\label{t:entry}
Write the entry payoff as $\Omega(K;A)=\Pi(\psi^{-1}(K/A))$. Under quantity
competition with $\eta=1$, with $a=2(\sigma-1)$, $b=\sigma-2$:
\[
\textbf{(A)}\ \Omega_{KK}<0:\quad
\varepsilon_G(S)=S\Bigl[\tfrac{a}{1+aS}-\tfrac{b}{1+bS}-\tfrac{\sigma}{1-S}\Bigr]<0
\ \text{at every }S\ (\sigma\ge2);
\]
\[
\textbf{(B)}\ \Omega_{KA}<0:\quad
\varepsilon_G+\varepsilon_\psi=\tfrac{aS}{1+aS}-\tfrac{bS}{1+bS}+\tfrac{1-2S}{1-S}>0
\ \text{at every }S\le\tfrac12,\ \text{with slack exactly }1/\sigma\text{ at }S=\tfrac12.
\]
If no group's share exceeds one half on the region between candidate outcomes (the
computed solution sits at $\approx0.30$ against an exact frontier of $0.545$), there
is \emph{at most one} entry outcome --- with no order of moves and no selection rule,
and with groups still competing with themselves.
\end{theorem}

\noindent\emph{Proof idea.} (A) makes each group's choice a cutoff on its own ladder
of factory sets; (B) makes that cutoff fall when the market toughens; and $A$ rises
when any group strengthens --- so two equilibria contradict each other. Brute force
over \emph{every} configuration: 120 of 120 random markets have exactly one outcome.
\src{The standard fixes fail here: cost-ranked entry (Gaubert--Itskhoki 2021) needs
every entrant separate, and an assumed order of moves (Yang 2023) makes market
structure an artifact of the order. The theorem replaces both.}

\paragraph{Wages, the last layer.} Input prices are a contraction (modulus
$\max_k\nu_k<1$); incomes are a linear system. For wages, let
$M_{nj}=\partial\ln W_n/\partial\ln w_j$ (unit row sums by homogeneity). The solver's
own map $\Psi_\kappa(w)_n=w_n(\mathrm{LD}_n/L_n)^\kappa$ has log-Jacobian
$(1-\kappa)I+\kappa M$; wherever that matrix is entrywise positive --- gross
substitutes off the diagonal, enough damping on it ($\kappa<1/(1-\min_nM_{nn})$, a
bound the model itself supplies) --- Birkhoff's (1957) theorem makes $\Psi_\kappa$ a
strict contraction in Hilbert's metric: \textbf{at most one equilibrium there, with a
geometric rate}. The positivity holds under two checked conditions (no group above
half a market --- the same one half as Theorem~\ref{t:entry} --- plus a demand
condition), holds at every point tested out to wage spreads of $e^2$, and is
\textbf{computer-certified on the continuum} of every wage vector within $\pm15\%$ of
the solution by outward-rounded interval arithmetic (Moore 1966) --- no gap between
grid points. What is \emph{not} available is an unconditional theorem, for a stated
reason: the one channel that can break the condition is the ordinary income effect,
which is this paper's own subject. A counterexample and the full honesty ledger are
in the guide. Ten dispersed starts reach the same wages \emph{and the same set of
operating firms}; every accounting identity holds to $10^{-15}$ away from equilibrium.

\section{Quantification: The Six Facts}

\begin{center}\footnotesize
\begin{tabular}{clccl}
\toprule
\# & Fact & Model & Data & Status\\
\midrule
1 & MNE export share; foreign dominant & $0.69$ ($0.54/0.15$) & $0.46$--$0.74$ & level fitted; \textbf{split produced}\\
2 & foreign firms in complex goods & $+0.42$ & $+0.43$ & fitted\\
3 & parents from few countries & $0.27$ & $0.13$ & \textbf{produced}, overshoots\\
4 & grouping by owner raises concentration & $\times1.94$ & $\times1.12$ & \textbf{produced} --- a prediction\\
5 & MNEs present $\Rightarrow$ local exports & $+46\%^{\dagger}$ & positive & \textbf{positive in the $\lambda$ config.}\\
6 & distance matters less for MNEs & $-0.52/{+}0.70$ & $-0.16/{+}0.05$ & \textbf{produced}, too large\\
\bottomrule
\end{tabular}
\end{center}

\noindent Three parameters are fitted (entry cost, head-office disadvantage, pool
scaling); Facts 3, 4 and 6 are genuine predictions.

\noindent\textbf{Fact 5, honestly.} The empirical fact survives a saturated
destination$\times$product$\times$year re-estimation at half its published size, so it
is real --- and the \emph{core} model gets it backwards ($-97\%$): with market
spending fixed, entry is zero-sum, and no cost-side spillover can fix that (measured,
not assumed). What fixes it is the missing margin: a \textbf{rest-of-world supplier in
every market}, so that sample exporters hold the small in-sample share
$\lambda\approx0.07$ seen in the data and business stealing falls on the outside
supply. $^{\dagger}$With that supplier plus two spillovers to single-plant locals
(productivity: Javorcik 2004; export infrastructure: Aitken--Hanson--Harrison 1997),
Fact~5 turns \textbf{positive on both margins} ($+46\%$; Fact 1 $=0.49$ with the split
intact, Fact 4 improves to $\times1.67$, one sign inside Fact 6 flips --- disclosed).
This \emph{six-fact configuration} (\texttt{julia mne\_model.jl lambda}) costs three
more fitted parameters and is reported beside the core, not in its place: six facts
against six fitted parameters is a different claim from five against three, and the
reader should see both.

\bigskip
\noindent\emph{Everything here is one file, base Julia, no packages
(\texttt{mne\_model.jl}); every number above has a committed run behind it. The full
document --- every proof written out, every choice justified where it is made, and a
provenance table naming the source of each component --- is \texttt{model\_guide.pdf}.}

\end{document}
